%% Lecture 10 -- Maxwell's equations completed; discrete symmetries; EM Lorentz invariants
\section{Lecture 10 -- Maxwell's Equations, Discrete Symmetries, and Lorentz Invariants of the EM Field}\label{sec:lec10}

\begin{center}
\begin{tabular}{ll}
  \textbf{Date:}\quad 10/05/2026 & \\
\end{tabular}
\end{center}
\hrule\vspace{1.5em}

\subsection*{Overview}

The previous lecture prepared all ingredients needed to derive Maxwell's equations from the action
principle.  The total action
\begin{equation}\label{eq:lec10:total-action}
  S = -\sum_n m_n c\int ds_n
  -\frac{1}{c}\sum_n e_n\int A_\mu\,dx_n^\mu
  -\frac{1}{16\pi c}\int F_{\mu\nu}F^{\mu\nu}\,d^4x
\end{equation}
was varied with respect to the field $A_\mu$, giving $\delta S_{\rm field}=(4\pi c)^{-1}\int(\partial_\mu F^{\mu\nu})\delta A_\nu\,d^4x$.
The four-current $J^\mu=(c\rho,\mathbf{J})$ was introduced and shown to be a four-vector obeying
$\partial_\mu J^\mu=0$.

This lecture:
\begin{itemize}
  \item completes the variational derivation of the inhomogeneous Maxwell equations;
  \item shows how the covariant form implies the continuity equation automatically;
  \item unpacks the covariant equation into the familiar 3-vector Maxwell equations;
  \item studies the behaviour of $\bE$ and $\bB$ under time-reversal ($t\to-t$) and spatial
        inversion ($\mathbf{r}\to-\mathbf{r}$), introducing polar and axial vectors;
  \item constructs the two independent Lorentz scalars built from $F_{\mu\nu}$ and extracts
        physical consequences about which configurations can be related by a boost.
\end{itemize}

%------------------------------------------------------------
\subsection{Completing the Derivation of Maxwell's Equations}\label{subsec:lec10-maxwell}

\subsubsection*{Physical Motivation}
The action has two types of degrees of freedom: particle worldlines and the field $A_\mu$.  We already
varied with respect to the particles and obtained the covariant Lorentz force.  Now we impose
stationarity under variations of the field.  Because $S_{\rm particle}$ does not depend on $A_\mu$,
only the last two terms contribute.

\subsubsection*{Variation of the interaction term}

Writing the interaction in terms of the four-current (derived in Lecture~09),
\begin{equation}\label{eq:lec10:Sint-J}
  S_{\rm int}
  =
  -\frac{1}{c^2}\int J^\mu A_\mu\,d^4x ,
\end{equation}
a shift $A_\mu\to A_\mu+\delta A_\mu$ (keeping particle trajectories fixed, so $J^\mu$ is unchanged)
gives
\begin{equation}\label{eq:lec10:dSint}
  \delta S_{\rm int}
  =
  -\frac{1}{c^2}\int J^\nu\,\delta A_\nu\,d^4x .
\end{equation}
Note: $J^\mu$ does not vary because it depends only on the particle degrees of freedom, which are
held fixed.

\subsubsection*{Combining the two variations}

Adding \eqref{eq:lec10:dSint} to the field variation derived in Lecture~09,
\begin{equation}\label{eq:lec10:dSfield}
  \delta S_{\rm field}
  =
  \frac{1}{4\pi c}\int(\partial_\mu F^{\mu\nu})\,\delta A_\nu\,d^4x ,
\end{equation}
the total variation is
\begin{equation}\label{eq:lec10:dS-total}
  \delta S
  =
  \int\!\left[
    \frac{1}{4\pi c}\,\partial_\mu F^{\mu\nu}
    -\frac{1}{c^2}\,J^\nu
  \right]\delta A_\nu\,d^4x .
\end{equation}
By the principle of least action, $\delta S=0$ for every admissible $\delta A_\nu$.  Since $\delta A_\nu$
is arbitrary, the bracket must vanish at every spacetime point:
\begin{equation}\label{eq:lec10:maxwell-covariant}
  \boxed{
  \partial_\mu F^{\mu\nu}
  =
  \frac{4\pi}{c}\,J^\nu .
  }
\end{equation}
\textit{Physical significance:} These are the \emph{inhomogeneous Maxwell equations} in compact covariant
form.  The left-hand side is determined entirely by the geometry of $A_\mu$; the right-hand side is the
four-current source.  The equation is manifestly Lorentz covariant: both sides are four-vectors.

\subsubsection*{The full set of Maxwell equations}

Together with the particle equation of motion (from varying the worldlines),
\begin{equation}\label{eq:lec10:particle-eom}
  \boxed{
  m\,\frac{du^\mu}{d\tau}
  =
  \frac{e}{c}\,F^\mu{}_\nu\,u^\nu ,
  }
\end{equation}
equations \eqref{eq:lec10:maxwell-covariant} and \eqref{eq:lec10:particle-eom} encode \emph{all of
classical electrodynamics} in just two compact lines.  This economy is possible because we understood
the system's symmetries well enough to write a natural action from first principles.

%------------------------------------------------------------
\subsection{Continuity Equation as an Automatic Consequence}\label{subsec:lec10-continuity}

\subsubsection*{Physical Motivation}
Charge conservation says that charge cannot be created or destroyed.  In a field theory derived from
an action, conservation laws should emerge automatically, either from symmetries (Noether) or from
algebraic identities.  Maxwell's equations contain charge conservation as an identity.

Apply $\partial_\nu$ to both sides of \eqref{eq:lec10:maxwell-covariant}:
\begin{equation}\label{eq:lec10:dnu-maxwell}
  \partial_\nu\partial_\mu F^{\mu\nu}
  =
  \frac{4\pi}{c}\,\partial_\nu J^\nu .
\end{equation}
The left-hand side vanishes identically.  The operator $\partial_\nu\partial_\mu$ is symmetric in
$\mu\leftrightarrow\nu$, while $F^{\mu\nu}$ is antisymmetric.  Contracting a symmetric operator with
an antisymmetric tensor gives zero:
\begin{equation}\label{eq:lec10:sym-anti}
  \partial_\nu\partial_\mu F^{\mu\nu}
  =
  -\partial_\nu\partial_\mu F^{\nu\mu}
  =
  -\partial_\mu\partial_\nu F^{\mu\nu}
  = 0 .
\end{equation}
(The second step uses $F^{\mu\nu}=-F^{\nu\mu}$; the third relabels dummy indices.)  Therefore
\begin{equation}\label{eq:lec10:continuity}
  \boxed{
  \partial_\mu J^\mu = 0 .
  }
\end{equation}
\textit{Physical significance:} Local charge conservation is not an \emph{additional} assumption; it
follows from Maxwell's equations automatically, as a consequence of the antisymmetry of $F^{\mu\nu}$.
This is an internal consistency check: the field equations can only be self-consistent if their source
is conserved.

\subsubsection*{Standard 3-vector derivation (cross-check)}

The same result follows from the familiar vector equations, providing a useful consistency check.
Take the divergence of the Ampère--Maxwell law \eqref{eq:lec10:ampere}:
\begin{equation}\label{eq:lec10:div-ampere}
  \nabla\cdot(\nabla\times\bB)
  - \frac{1}{c}\frac{\partial}{\partial t}(\nabla\cdot\bE)
  =
  \frac{4\pi}{c}\,\nabla\cdot\\Jvec .
\end{equation}
The first term vanishes because the divergence of any curl is zero.  Substituting Gauss's law
\eqref{eq:lec10:gauss}, $\nabla\cdot\bE=4\pi\rho$:
\begin{equation}\label{eq:lec10:continuity-3vec}
  0 - \frac{4\pi}{c}\frac{\partial\rho}{\partial t}
  =
  \frac{4\pi}{c}\,\nabla\cdot\\Jvec,
\end{equation}
which rearranges to
\begin{equation}\label{eq:lec10:continuity-3vec-final}
  \boxed{
  \frac{\partial\rho}{\partial t} + \nabla\cdot\\Jvec = 0 .
  }
\end{equation}
\textit{Physical significance:} Maxwell's equations are only mathematically self-consistent if charge
is conserved.  The displacement current term $\partial_t\bE/c$ that Maxwell added to Ampère's law is
precisely what ensures this consistency.

%------------------------------------------------------------
\subsection{Explicit Component Form of Maxwell's Equations}\label{subsec:lec10-components}

\subsubsection*{Physical Motivation}
The covariant equation \eqref{eq:lec10:maxwell-covariant} contains four equations (one for each value
of $\nu$).  Let us unpack them.  Recall (from Lecture~08) the explicit matrix of $F^{\mu\nu}$ in
Gaussian CGS:
\begin{equation}\label{eq:lec10:F-matrix}
  F^{\mu\nu} =
  \begin{pmatrix}
    0 & -E_x & -E_y & -E_z \\
    E_x & 0 & -B_z & B_y \\
    E_y & B_z & 0 & -B_x \\
    E_z & -B_y & B_x & 0
  \end{pmatrix},
\end{equation}
so $F^{0i}=-E_i$ (and $F^{i0}=E_i$), and $F^{ij}=-\epsilon^{ijk}B_k$.

\subsubsection*{The $\nu=0$ component: Gauss's law}

\begin{align}
  \partial_\mu F^{\mu 0}
  &=
  \partial_i F^{i0}
  =
  \partial_i E_i
  =
  \nabla\cdot\bE ,
  \label{eq:lec10:gauss-lhs}
\end{align}
where we used $F^{00}=0$.  The right-hand side is $(4\pi/c)\,J^0 = (4\pi/c)(c\rho)=4\pi\rho$.
Therefore
\begin{equation}\label{eq:lec10:gauss}
  \boxed{
  \nabla\cdot\bE = 4\pi\rho .
  }
\end{equation}
\textit{Physical significance:} This is Gauss's law.  Electric field lines originate and terminate on
charge.

\subsubsection*{The $\nu=1$ component: Ampère--Maxwell law ($x$-direction)}

\begin{align}
  \partial_\mu F^{\mu 1}
  &=
  \partial_0 F^{01}
  +\partial_2 F^{21}
  +\partial_3 F^{31}
  \notag\\
  &=
  -\frac{1}{c}\partial_t E_x
  +\partial_y B_z
  -\partial_z B_y
  \notag\\
  &=
  (\nabla\times\bB)_x - \frac{1}{c}\frac{\partial E_x}{\partial t},
  \label{eq:lec10:ampere-x-lhs}
\end{align}
where we computed:
$F^{01}=-E_x$ so $\partial_0 F^{01} = -(1/c)\partial_t E_x$;
$F^{21}=-F^{12}=B_z$ so $\partial_2 F^{21}=\partial_y B_z$;
$F^{31}=-F^{13}=-B_y$ so $\partial_3 F^{31}=-\partial_z B_y$.
The right-hand side is $(4\pi/c)J^1=(4\pi/c)J_x$.  Repeating for $\nu=2,3$ gives the full equation
\begin{equation}\label{eq:lec10:ampere}
  \boxed{
  \nabla\times\bB - \frac{1}{c}\frac{\partial\bE}{\partial t}
  =
  \frac{4\pi}{c}\,\Jvec .
  }
\end{equation}
\textit{Physical significance:} Currents and changing electric fields source the magnetic field.  The
term $(1/c)\partial_t\bE$ is Maxwell's displacement current; without it, $\nabla\times\bB$ would not
be divergence-free, violating $\partial_\mu J^\mu=0$.

\subsubsection*{The homogeneous equations: direct derivation from definitions}

\subsubsection*{Physical Motivation}
The fundamental object in our theory is the four-potential $A^\mu=(\phi,\Avec)$.  All physical fields
are built from it.  The two definitions
\begin{align}
  \bE &= -\nabla\phi - \frac{1}{c}\frac{\partial\Avec}{\partial t}, \label{eq:lec10:E-def}\\
  \bB &= \nabla\times\Avec, \label{eq:lec10:B-def}
\end{align}
are \emph{not} dynamical equations; they are kinematic identities that encode the choice to describe
electromagnetism via a potential.  The homogeneous Maxwell equations are then \emph{automatic
consequences of these definitions}, requiring no additional input from the action.

\paragraph{Faraday's law.}
Take the curl of both sides of \eqref{eq:lec10:E-def}:
\begin{equation}\label{eq:lec10:faraday-start}
  \nabla\times\bE
  =
  -\nabla\times(\nabla\phi)
  -\frac{1}{c}\frac{\partial}{\partial t}(\nabla\times\Avec).
\end{equation}
The first term vanishes identically: the curl of any gradient is zero ($\nabla\times\nabla\phi=0$).
In the second term, partial derivatives with respect to space and time commute, and
$\nabla\times\Avec=\bB$ by definition \eqref{eq:lec10:B-def}.  Therefore
\begin{equation}\label{eq:lec10:faraday}
  \boxed{
  \nabla\times\bE = -\frac{1}{c}\frac{\partial\bB}{\partial t}.
  }
\end{equation}
\textit{Physical significance:} This is Faraday--Lenz's law.  A time-varying magnetic field induces
a curling electric field.  It is a pure consequence of the potential structure, not of any dynamics.

\paragraph{No magnetic monopoles ($\nabla\cdot\bB=0$).}
Take the divergence of \eqref{eq:lec10:B-def}:
\begin{equation}\label{eq:lec10:divB-calc}
  \nabla\cdot\bB
  =
  \nabla\cdot(\nabla\times\Avec) = 0,
\end{equation}
since the divergence of any curl vanishes identically.  Therefore
\begin{equation}\label{eq:lec10:divB}
  \boxed{
  \nabla\cdot\bB = 0 .
  }
\end{equation}
\textit{Physical significance:} There are no magnetic charges (monopoles) in this theory.  The result
is not an experimental statement but a \emph{definition}: we chose $\bB=\nabla\times\Avec$, and for
any vector field that is a curl, the divergence is zero.

\paragraph{Covariant statement (Bianchi identity).}
Both results are unified in the covariant identity
\begin{equation}\label{eq:lec10:bianchi}
  \partial_\lambda F_{\mu\nu}
  +\partial_\mu F_{\nu\lambda}
  +\partial_\nu F_{\lambda\mu}
  = 0 ,
\end{equation}
which holds for any $F_{\mu\nu}=\partial_\mu A_\nu-\partial_\nu A_\mu$.  Taking $\lambda=0$ recovers
Faraday's law; contracting over the spatial indices recovers $\nabla\cdot\bB=0$.

\subsubsection*{The E–B asymmetry in the definitions}

Before asking about monopoles, notice that the two definitions
\eqref{eq:lec10:E-def}--\eqref{eq:lec10:B-def} are \emph{not symmetric} in $\bE$ and $\bB$:
\begin{align*}
  \bE &= -\nabla\phi - \frac{1}{c}\frac{\partial\Avec}{\partial t}
        \quad (\text{gradient term} + \text{time derivative of }\Avec), \\
  \bB &= \nabla\times\Avec
        \quad (\text{curl of }\Avec\text{ only}).
\end{align*}
There is no symmetry between the electric and magnetic sectors at the level of the definitions,
because the action contains only \emph{one} four-potential $A^\mu$ and only \emph{electric} charges.
This asymmetry is a direct reflection of the experimental fact that isolated magnetic charges have
never been observed.

\subsubsection*{Why $\nabla\cdot\bB=0$ is not a fundamental law of nature}

The statement $\nabla\cdot\bB=0$ is a consequence of our choice to describe the field by the
potential $A^\mu$.  If magnetic monopoles existed, the force law
$\mathbf{F}=e(\bE+\mathbf{v}/c\times\bB)$ would need an additional magnetic-charge term.  We have
no such term in our action because we have never observed magnetic monopoles in the laboratory.

If monopoles were discovered, one would introduce a \emph{second} four-potential
$A^\mu_{\rm mag}=(\phi_M,\Avec_M)$, and the definitions would be extended to
\begin{align}
  \bE &= -\nabla\phi - \frac{1}{c}\dot{\Avec} + \nabla\times\Avec_M , \label{eq:lec10:E-monopole}\\
  \bB &= \nabla\times\Avec - \nabla\phi_M - \frac{1}{c}\dot{\Avec}_M , \label{eq:lec10:B-monopole}
\end{align}
making the electric and magnetic sectors fully symmetric: $\bE$ acquires a curl of $\Avec_M$ and
$\bB$ acquires the same gradient + time-derivative structure as $\bE$.  The equation
$\nabla\cdot\bB=0$ would be replaced by
\begin{equation}\label{eq:lec10:divB-monopole}
  \nabla\cdot\bB = 4\pi\rho_M .
\end{equation}

\subsubsection*{Why we do not simply add monopoles speculatively}

Introducing a monopole potential immediately predicts a new magnetic force.  If that force were
comparable in strength to the electric force, it would have been detected long ago.  Any consistent
monopole theory must therefore contain a new coupling constant that suppresses the monopole
interaction to below current experimental sensitivity.

An analogy: general relativity predicts gravitational waves, but they are extraordinarily weak
because Newton's constant $G$ is very small.  If one had never detected gravity directly, one would
still have to include $G$ as a free parameter that the theory cannot fix on symmetry grounds alone.
Similarly, a monopole theory would require a new coupling constant to explain the absence of observed
magnetic-charge effects.

\textit{Physical significance:} The framework of action + symmetry + potential does not forbid
monopoles on logical grounds.  Their absence is an experimental input, encoded in the choice of which
four-potential we include in the action.

%------------------------------------------------------------
\subsection{The Methodology of Analytical Field Theory}\label{subsec:lec10-methodology}

\subsubsection*{Physical Motivation}
The derivation of Maxwell's equations from the action is not just a mathematical trick.  It reflects
a deep methodological principle that underlies all of modern theoretical physics.  The professor
describes it as \emph{analytical mechanics of field theory}: the same philosophy that Lagrange and
Hamilton applied to mechanics, now applied to fields.

\subsubsection*{Analytical mechanics as the prototype}

In Newtonian mechanics one writes forces and integrates $F=ma$.  Analytical mechanics
(Lagrangian/Hamiltonian) produces \emph{exactly the same equations} but arrives at them from two
ingredients:
\begin{enumerate}
  \item \textbf{Symmetry principle:} the action must be invariant under the symmetries of the system
        (Galilean invariance for non-relativistic mechanics).
  \item \textbf{Principle of least action:} the physical trajectory extremises $S=\int L\,dt$.
\end{enumerate}
The Lagrangian itself is not derived from first principles; one must know from experiment what
interactions are present (e.g.\ a $1/r$ potential for the Kepler problem).  The symmetry + action
framework then generates the equations of motion unambiguously.

\subsubsection*{The same logic applied to electrodynamics}

This course performs exactly the same programme for the electromagnetic field:
\begin{enumerate}
  \item \textbf{Symmetry:} Lorentz invariance (special relativity), gauge invariance, and the
        superposition principle.
  \item \textbf{Experimental inputs:} The observed force on a charge is $\mathbf{F}=e(\bE+\mathbf{v}/c\times\bB)$
        (Lorentz force), not a magnetic-monopole force.  The superposition principle is observed.
  \item \textbf{Principle of least action:} the physical fields and trajectories extremise
        \eqref{eq:lec10:total-action}.
\end{enumerate}
Together these three inputs \emph{uniquely} determine the action and hence all of Maxwell's
equations.  As the professor put it in Lecture~01: ``this is analytical mechanics of field theory.''

\subsubsection*{Why experiment is irreplaceable}

Symmetry constraints are necessary but not sufficient to fix a theory.  There are in general many
actions compatible with a given symmetry group.  Experimental observations select which terms appear
and fix the coupling constants.

\begin{itemize}
  \item In the Kepler problem, symmetry allows any central force; only experiment (Kepler's laws)
        selects $F\propto 1/r^2$.
  \item In electrodynamics, Lorentz + gauge invariance allows both $F_{\mu\nu}F^{\mu\nu}$ and
        $\epsilon^{\mu\nu\alpha\beta}F_{\mu\nu}F_{\alpha\beta}$ in the Lagrangian; the latter is a
        total derivative and is discarded because it contributes nothing to the equations of motion.
        The coupling $e/c$ in the interaction term is fixed by measuring the electron charge.
\end{itemize}

\subsubsection*{The Standard Model as the same programme applied further}

The same methodology underlies the Standard Model of particle physics, built roughly 50--60 years
ago:
\begin{enumerate}
  \item Identify the symmetry group (Lorentz invariance + gauge symmetries $SU(3)\times SU(2)\times
        U(1)$).
  \item Write all actions (Lagrangian terms) consistent with those symmetries.
  \item Use experimental data (particle masses, decay rates, scattering cross-sections) to fix the
        free parameters and eliminate terms not seen in nature.
\end{enumerate}
The result is a compact set of equations that encode all known non-gravitational physics.  The guiding
principle throughout is: \emph{action + symmetry, then constrain with experiment}.

\textit{Physical significance:} Electrodynamics, studied in this course, is both the simplest and
the historical prototype of this programme.  Mastering it is the prerequisite for understanding how
all fundamental interactions are formulated.

%------------------------------------------------------------
\subsection{Aside: Feynman's Path-Integral Formulation}\label{subsec:lec10-path-integral}

\subsubsection*{Physical Motivation}
The discussion of the action principle prompted a brief digression on Feynman's alternative
formulation of quantum mechanics, which generalises the least-action idea from classical to quantum
physics.

In classical mechanics, a particle follows the single trajectory that \emph{extremises} the action
$S$: of all conceivable paths from $(\mathbf{x}_1,t_1)$ to $(\mathbf{x}_2,t_2)$, only the one
satisfying the Euler--Lagrange equations is realised.

Feynman proposed that in quantum mechanics every path is realised \emph{simultaneously}.  The
probability amplitude for propagation from point 1 to point 2 is the sum (integral) over all paths,
each weighted by the phase factor
\begin{equation}\label{eq:lec10:path-integral}
  \boxed{
  \mathcal{A}_{1\to 2}
  \propto
  \int_{\text{all paths}} \mathcal{D}[x(t)]\;
  e^{iS[x(t)]/\hbar} ,
  }
\end{equation}
where $S[x(t)]$ is the classical action evaluated along a given path and $\mathcal{D}[x(t)]$ denotes
integration over the space of all trajectories.

\textit{Physical significance:} In the classical limit $\hbar\to 0$, the phase $e^{iS/\hbar}$
oscillates infinitely rapidly for paths far from the classical one, so they cancel by destructive
interference.  Only the neighbourhood of the \emph{stationary-phase path} (i.e.\ the classical
trajectory) contributes, recovering the principle of least action.  Quantum mechanics is therefore
the deeper framework: classical mechanics is its stationary-phase approximation.

%------------------------------------------------------------
\subsection{Discrete Symmetries: Time Reversal and Spatial Inversion}\label{subsec:lec10-discrete}

\subsubsection*{Physical Motivation}
The proper orthochronous Lorentz group $SO^+(3,1)$ is continuously connected to the identity.  But
there are also discrete transformations: time reversal $T:t\to-t$ and spatial parity $P:\mathbf{r}\to
-\mathbf{r}$ (in three dimensions).  Electromagnetism \emph{is} symmetric under both, but in nature
the weak force violates $T$ and $P$ (and even $CP$).  Here we determine how $\bE$ and $\bB$ transform
under these discrete symmetries by using Maxwell's equations as constraints.

\subsubsection*{Time reversal $t\to-t$}

\textbf{Charge density:} $\rho(\mathbf{x},t)\to\rho(\mathbf{x},-t)$.  A static charge is unaffected;
the value of $\rho$ does not change sign:
\begin{equation}\label{eq:lec10:T-rho}
  \rho \;\longrightarrow\; \rho .
\end{equation}
\textbf{Current density:} $\mathbf{J}=\rho\mathbf{v}$.  The velocity $\mathbf{v}=d\mathbf{x}/dt$
changes sign under $t\to-t$, so
\begin{equation}\label{eq:lec10:T-J}
  \mathbf{J} \;\longrightarrow\; -\mathbf{J} .
\end{equation}
\textbf{Electric field:} Gauss's law $\nabla\cdot\bE=4\pi\rho$ relates $\bE$ to $\rho$.  The right-hand
side does not change sign.  The spatial derivatives $\partial_i$ are unaffected by $t\to-t$.  Therefore
\begin{equation}\label{eq:lec10:T-E}
  \bE \;\longrightarrow\; \bE .
\end{equation}
\textbf{Magnetic field:} Consider Faraday's law $\nabla\times\bE=-(1/c)\partial_t\bB$.  The left-hand
side transforms as $\bE\to\bE$ (purely spatial curl), so it is unchanged.  The right-hand side
acquires a sign from $\partial_t\to-\partial_t$ and another from $\bB\to\;?$\,.  For the equation
to remain satisfied, the two signs must cancel:
\begin{equation}\label{eq:lec10:T-B}
  \bB \;\longrightarrow\; -\bB .
\end{equation}

\begin{center}
\begin{tikzpicture}[>=Latex, scale=1.0]
  \node[align=center] at (-2.2,0) {$t\to -t$};
  \draw[->] (-1.4,0.4) -- (-0.3,0.9) node[midway,above,font=\small] {$\bE$};
  \draw[->] (-1.4,-0.4) -- (-0.3,-0.9) node[midway,below,font=\small] {$\bB$};
  \node[right,font=\small] at (-0.2,0.9) {unchanged};
  \node[right,font=\small] at (-0.2,-0.9) {flips sign};
\end{tikzpicture}
\end{center}

\subsubsection*{Spatial inversion $\mathbf{r}\to-\mathbf{r}$}

\textbf{Charge and current:} $\rho$ is a scalar density; it is unchanged.  $\mathbf{J}=\rho\mathbf{v}$
changes sign because $\mathbf{v}=d\mathbf{x}/dt$ changes sign:
\begin{equation}\label{eq:lec10:P-rho-J}
  \rho \;\longrightarrow\; \rho, \qquad \mathbf{J} \;\longrightarrow\; -\mathbf{J} .
\end{equation}
\textbf{Electric field:} Under $\mathbf{r}\to-\mathbf{r}$, the spatial gradient $\nabla_i\to-\nabla_i$.
Gauss's law $\nabla\cdot\bE=4\pi\rho$ then requires $\bE\to-\bE$ (to cancel the sign of
$\nabla\to-\nabla$):
\begin{equation}\label{eq:lec10:P-E}
  \bE \;\longrightarrow\; -\bE .
\end{equation}
\textbf{Magnetic field:} In Faraday's law, the curl $\nabla\times\bE\to-(-\bE)=+(\nabla\times\bE)$
(two sign flips cancel).  The time derivative $\partial_t$ is unaffected.  Therefore $\bB$ must be
unchanged to keep the equation satisfied:
\begin{equation}\label{eq:lec10:P-B}
  \bB \;\longrightarrow\; +\bB .
\end{equation}

\begin{center}
\renewcommand{\arraystretch}{1.3}
\begin{tabular}{ccc}
  \toprule
  & Time reversal ($t\to-t$) & Spatial inversion ($\mathbf{r}\to-\mathbf{r}$) \\
  \midrule
  $\rho$ & $+\rho$ & $+\rho$ \\
  $\mathbf{J}$ & $-\mathbf{J}$ & $-\mathbf{J}$ \\
  $\bE$ & $+\bE$ & $-\bE$ \\
  $\bB$ & $-\bB$ & $+\bB$ \\
  \bottomrule
\end{tabular}
\end{center}

\subsubsection*{Polar and axial vectors}

A \textbf{polar vector} (ordinary vector) changes sign under spatial inversion, just like $\mathbf{r}$
itself.  Examples: $\mathbf{r}$, $\bE$, $\mathbf{J}$.

An \textbf{axial vector} (pseudovector) does \emph{not} change sign under spatial inversion.  It
behaves like a rotation axis.  Examples: $\bB$, angular momentum $\mathbf{L}=\mathbf{r}\times\mathbf{p}$.

\textit{Physical significance:} $\bB$ is axial because it is built as a curl (a cross product), which
inherits an extra sign flip that compensates the inversion.  This distinction matters when writing
Lorentz-invariant combinations: only contractions that respect the parity of each vector produce true
scalars.

\subsubsection*{Caveat: parity in two dimensions}

A common error: in two spatial dimensions, $\mathbf{r}\to-\mathbf{r}$ is a rotation by $180°$, not a
parity transformation, because $\det(-\mathbf{1}_{2\times 2})=+1$.  True parity requires an odd number
of sign flips, i.e.\ $\det(\text{transformation})=-1$.  In three dimensions, $\mathbf{r}\to-\mathbf{r}$
has $\det=-1$ and is a genuine parity.

%------------------------------------------------------------
\subsection{Lorentz Invariants of the Electromagnetic Field}\label{subsec:lec10-invariants}

\subsubsection*{Physical Motivation}
The components $\bE$ and $\bB$ mix under Lorentz boosts (they are entries of the rank-2 tensor
$F^{\mu\nu}$, not separate four-vectors).  Yet certain combinations must be frame-independent.  These
Lorentz scalars encode absolute properties of the electromagnetic field: which configurations are
related by a boost, and which are not.

\subsubsection*{First invariant: $F_{\mu\nu}F^{\mu\nu}$}

The free-field Lagrangian density $\mathcal{L}=-F_{\mu\nu}F^{\mu\nu}/(16\pi)$ is a Lorentz scalar
(it enters the Lorentz-scalar action).  Computing $F_{\mu\nu}F^{\mu\nu}$ explicitly:

The matrix $F_{\mu\nu}$ (indices both down) differs from $F^{\mu\nu}$ by raising both indices, i.e.\
multiplying by $\eta^{\mu\alpha}\eta^{\nu\beta}$.  Raising a timelike index (0) costs no sign;
raising a spacelike index costs a $-1$ factor per index.  Therefore $F^{i0}=-F_{i0}$ and
$F^{ij}=F_{ij}$ (both indices spatial, two minus signs cancel).  The time-space components thus pick
up a relative sign:
\begin{align}
  F_{\mu\nu}F^{\mu\nu}
  &=
  2\!\left[F_{01}F^{01}+F_{02}F^{02}+F_{03}F^{03}\right]
  +
  2\!\left[F_{12}F^{12}+F_{13}F^{13}+F_{23}F^{23}\right]
  \notag\\
  &=
  2(-E_x^2-E_y^2-E_z^2)
  +
  2(B_z^2+B_y^2+B_x^2)
  \notag\\
  &=
  2(B^2-E^2).
  \label{eq:lec10:I1-calc}
\end{align}
(Each off-diagonal pair contributes twice because of the symmetry of the sum.)
\begin{equation}\label{eq:lec10:I1}
  \boxed{
  F_{\mu\nu}F^{\mu\nu}
  =
  2\!\left(B^2 - E^2\right).
  }
\end{equation}
\textit{Physical significance:} The sign of $B^2-E^2$ is a Lorentz-invariant property.  An
electromagnetic configuration with $B^2>E^2$ remains so in every frame; a pure-electric configuration
($B=0$, $E\neq 0$) has $B^2-E^2<0$ and stays that way.  In particular, it is impossible to boost to
a frame where $B=0$ if $B^2-E^2>0$ in the original frame.

\subsubsection*{Second invariant: the Levi-Civita pseudoscalar}

There is a second quadratic combination,
\begin{equation}\label{eq:lec10:I2-def}
  \epsilon^{\alpha\beta\mu\nu}\,F_{\alpha\beta}\,F_{\mu\nu},
\end{equation}
where $\epsilon^{\alpha\beta\mu\nu}$ is the \textbf{four-dimensional Levi-Civita symbol}:
\begin{equation}\label{eq:lec10:eps-def}
  \epsilon^{0123}=+1, \qquad
  \text{completely antisymmetric in all four indices.}
\end{equation}
This symbol is a \emph{pseudotensor} (it picks up a $\det\Lambda$ factor under Lorentz
transformations), so \eqref{eq:lec10:I2-def} is a pseudoscalar.

A direct calculation (summing over all index permutations) yields contributions of the form
$F_{0i}F_{jk}$, which combine to give $E_x B_x + E_y B_y + E_z B_z$:
\begin{equation}\label{eq:lec10:I2}
  \boxed{
  \epsilon^{\alpha\beta\mu\nu}\,F_{\alpha\beta}\,F_{\mu\nu}
  =
  -8\,\bE\cdot\bB .
  }
\end{equation}
\textit{Physical significance:} The angle between $\bE$ and $\bB$ (or rather, the sign of
$\bE\cdot\bB$) is a Lorentz-invariant property.  If $\bE\perp\bB$ in one frame, they are orthogonal
in all frames.

\paragraph{Why this term does not affect the field equations.}
At first sight, one might add $\alpha\,\epsilon^{\alpha\beta\mu\nu}F_{\alpha\beta}F_{\mu\nu}$ to the
Lagrangian as a second quadratic invariant.  However, one can show that it is a \emph{total
spacetime derivative}:
\begin{equation}\label{eq:lec10:I2-total-deriv}
  \epsilon^{\alpha\beta\mu\nu}F_{\alpha\beta}F_{\mu\nu}
  =
  \partial_\mu\!\left(\epsilon^{\mu\nu\alpha\beta}A_\nu F_{\alpha\beta}\right).
\end{equation}
Adding a total derivative to the Lagrangian density changes the action by a boundary term, which does
not affect the Euler--Lagrange equations.  This is why we discarded it when constructing the action
in Lecture~08.

%------------------------------------------------------------
\subsection{Physical Consequences of the Two Invariants}\label{subsec:lec10-consequences}

\subsubsection*{Physical Motivation}
The two invariants $B^2-E^2$ and $\bE\cdot\bB$ are the complete set of independent Lorentz scalars
quadratic in $F$.  They classify electromagnetic field configurations and determine what can be
transformed into what by a Lorentz boost.

\subsubsection*{Case 1: Radiation ($B=E$, $\bE\perp\bB$)}
For a plane electromagnetic wave (radiation), $|\bE|=|\bB|$ and $\bE\perp\bB$.  Therefore both
invariants vanish: $B^2-E^2=0$ and $\bE\cdot\bB=0$.  Since both invariants are zero, they remain
zero in every frame.  \emph{It is impossible to eliminate electromagnetic radiation by switching to
a different inertial frame.}  This is physically obvious: if you are ``burned'' by radiation in one
frame, you are burned in all frames.

\subsubsection*{Case 2: $B^2>E^2$}
One can always boost to a frame where $\bE'=0$ (purely magnetic field).  The invariant
$B^2-E^2>0$ is preserved and now reads $B'^2>0$, consistent.  But one \emph{cannot} boost to a frame
with $\bB'=0$, because that would require $B'^2-E'^2=-E'^2\leq 0$, contradicting $B^2-E^2>0$.

\subsubsection*{Case 3: $E^2>B^2$}
One can always boost to a frame where $\bB'=0$ (purely electric field).  Conversely, one cannot
boost to a frame with $\bE'=0$.

\subsubsection*{Case 4: $F^{\mu\nu}=0$}
If all components of $F^{\mu\nu}$ vanish in one frame, both invariants are zero and the tensor
remains zero in every frame (it transforms tensorially; $\Lambda F\Lambda^T=0$).

\begin{center}
\renewcommand{\arraystretch}{1.3}
\begin{tabular}{ccc}
  \toprule
  Condition (any frame) & Can boost to & Cannot boost to \\
  \midrule
  $B^2-E^2>0$ & $\bE'=0$ & $\bB'=0$ \\
  $B^2-E^2<0$ & $\bB'=0$ & $\bE'=0$ \\
  $B^2-E^2=0$, $\bE\cdot\bB=0$ (radiation) & neither & neither \\
  $\bE\cdot\bB=0$ & $\bE'\perp\bB'$ & $\bE'\parallel\bB'\neq 0$ \\
  \bottomrule
\end{tabular}
\end{center}

\textit{Physical significance:} The two invariants are the complete set of information that survives
a boost.  All other properties of $\bE$ and $\bB$ (individual magnitudes, directions relative to
coordinate axes) can be changed by choosing a different frame.

\subsection*{Summary}

\begin{itemize}
  \item Varying the total action with respect to $A_\mu$ gives the inhomogeneous Maxwell equations
        $\partial_\mu F^{\mu\nu}=(4\pi/c)J^\nu$.
  \item The covariant Lorentz force law $m\,du^\mu/d\tau = (e/c)F^\mu{}_\nu u^\nu$ and the field
        equation together encode all of classical electrodynamics.
  \item Applying $\partial_\nu$ to the field equation and using antisymmetry of $F^{\mu\nu}$
        immediately yields charge conservation $\partial_\mu J^\mu=0$.  The same result follows from
        the 3-vector equations by taking $\nabla\cdot$ of the Ampère--Maxwell law and substituting
        Gauss's law; Maxwell's displacement current is essential for consistency.
  \item The $\nu=0$ component gives Gauss's law; the spatial components give the Ampère--Maxwell law.
  \item The homogeneous pair ($\nabla\times\bE=-(1/c)\dot{\bB}$ and $\nabla\cdot\bB=0$) follows
        directly from the definitions $\bE=-\nabla\phi-(1/c)\dot{\Avec}$ and $\bB=\nabla\times\Avec$,
        via $\nabla\times\nabla f=0$ and $\nabla\cdot(\nabla\times\mathbf{F})=0$.  These are kinematic
        identities, not dynamical equations.  The absence of monopoles is an experimental input, not
        a logical necessity.
  \item The definitions of $\bE$ and $\bB$ in terms of $A^\mu$ are \emph{not} E--B symmetric; this
        asymmetry reflects the absence of magnetic charges.  Adding a second four-potential
        $A^\mu_{\rm mag}$ restores full symmetry and gives $\nabla\cdot\bB=4\pi\rho_M$.  Any
        monopole theory requires a new suppression constant to explain non-observation.
  \item The derivation exemplifies \emph{analytical field theory}: symmetry (Lorentz + gauge +
        superposition) + principle of least action + experimental inputs (Lorentz force, no monopoles)
        uniquely fix the action and hence all of Maxwell's equations.  This is the same programme
        used to construct the Standard Model.
  \item Under time reversal: $\bE\to+\bE$, $\bB\to-\bB$.  Under spatial inversion: $\bE\to-\bE$
        (polar vector), $\bB\to+\bB$ (axial vector).
  \item The first Lorentz invariant $F_{\mu\nu}F^{\mu\nu}=2(B^2-E^2)$ determines which of $\bE$,
        $\bB$ can be boosted away.
  \item The second Lorentz invariant $\epsilon^{\alpha\beta\mu\nu}F_{\alpha\beta}F_{\mu\nu}=-8\bE\cdot\bB$
        determines whether $\bE$ and $\bB$ are orthogonal in all frames.
  \item Radiation has both invariants zero; it cannot be eliminated by any boost.
\end{itemize}
